\documentclass{scrartcl}

\usepackage[utf8]{inputenc}
\usepackage[english]{babel}
\usepackage{amsmath}
\usepackage{amsfonts}
\usepackage{amssymb}
\usepackage{amsthm}
\usepackage{tikz}
\usepackage{hyphenat}
\usepackage{graphicx}
\usepackage{nicefrac}
\usepackage{booktabs}
\usepackage{todonotes}
\usepackage{bm}
\usepackage{BOONDOX-calo} % for nicer calligraphic symbols
\usepackage{authblk}
\usepackage{subcaption}

\usepackage[square]{natbib}
\bibliographystyle{chicago}

\usepackage{geometry}
\usepackage[
    unicode = true, 
    bookmarks = false, 
    breaklinks = true, 
    pdfborder = {0 0 1},
    colorlinks = true
    ]{hyperref}

\geometry{vmargin = 3cm, hmargin = 3cm}

\usetikzlibrary{
    automata, 
    positioning, 
    calc, 
    shapes, 
    arrows, 
    fit
    }

% \tikzset{
%     every loop/.style = {
%         min distance = 1mm, 
%         out = 25, 
%         in = -25, 
%         looseness = 5
%         }
%     }

\definecolor{PastelGreen}{HTML}{6ECB63}
\definecolor{EarthyGreen}{HTML}{588157}
\definecolor{Azure}{RGB}{0, 128, 255}
\definecolor{Fiery}{HTML}{c1121f}

\theoremstyle{definition}
\newtheorem{theorem}{Theorem}
\newtheorem{lemma}{Lemma}
\newtheorem{corollary}{Corollary}
\newtheorem{example}{Example}
\newtheorem{conjecture}{Conjecture}
\newtheorem{proposition}{Proposition}
\newtheorem{definition}{Definition}
\newtheorem{remark}{Remark}
\newtheorem{claim}{Claim}
\newtheorem{observation}{Observation}

\DeclareMathOperator*{\argmax}{arg\,max}
\DeclareMathOperator*{\argmin}{arg\,min}

\renewcommand{\descriptionlabel}[1]{$\triangleright$ \textit{#1}}

\renewcommand{\phi}{\varphi}
\renewcommand{\epsilon}{\varepsilon}
\renewcommand{\theta}{\vartheta}
\newcommand{\len}{\mathsf{len}}
\newcommand{\net}{\mathsf{net}}
\renewcommand{\min}{\mathit{min}}
\renewcommand{\max}{\mathit{max}}
\newcommand{\avg}{\mathit{avg}}
\newcommand{\ev}{\mathbb{E}}
\newcommand{\merge}{M}

\hypersetup{
    linkcolor = blue,
    citecolor = Azure,
    filecolor = blue,
    urlcolor = blue
    }

\tikzset{
	a/.style={
        draw,
        rounded corners=2pt, 
        minimum width=1cm, 
        minimum height=1cm, 
        thick,
        color=black,
        fill=EarthyGreen, 
        text=white,
        font=\Large
        },
    b/.style={
        draw,
        rounded corners=2pt, 
        minimum width=1cm, 
        minimum height=1cm, 
        thick,
        color=black,
        fill=Fiery, 
        text=white,
        font=\Large
        },
    a-agent/.style={
        draw,
        circle, 
        minimum width=2.5em, 
        thick,
        color=black,
        fill=EarthyGreen, 
        text=white,
        font=\Large
        },
    b-agent/.style={
        draw,
        circle, 
        minimum width=2.5em, 
        thick,
        color=black,
        fill=Fiery, 
        text=white,
        font=\Large
        },
    neutral-agent/.style={
        draw,
        circle, 
        minimum width=2.5em, 
        thick,
        color=black,
        fill=gray, 
        text=white,
        font=\Large
        }
    }


\newcommand{\adrian}[1]{
    {
        \hfill
        \begin{minipage}{0.9\textwidth}
            \begin{flushright}
                \color{PastelGreen}
                \footnotesize
                \small 
                {\footnotesize\textsf{\textbf{Adrian}}}\\ 
                \textsf{#1}
            \end{flushright}
        \end{minipage}
    }
    }


\title{
    Social Learning with Logical Formulas
    }
\author{Edoardo, Adrian \dots}
\date{}




\begin{document}
\maketitle
% \setcounter{tocdepth}{1}
% \tableofcontents

\begin{abstract}
    \noindent
    We study social learning with logical formulas.
\end{abstract}

\section{Introduction}\label{sec:introduction}
Does social learning work when opinions are more sophisticated representations of the world,
like logical formulas? We find out.

There is work that look at opinion dynamics with logical formulas
\citep{Schwind2015BeliefRevisiona,Schwind2016PromotingBeliefs,Christoff2017StabilityBinary,Botan2019MultiIssueOpinion,Schwind2024BeliefFlowFramework,Schwind2026CharacterizingConsensuses},
but usually with an eye towards consensus.
It seems that the truth-tracking dimension on a network has not been tackled yet.



\paragraph{Contributions.}
We show interesting things.

\paragraph{Related Literature.}
Belief merging operators \citep{Konieczny2011LogicBased} and their truth-tracking properties \citep{Everaere2010EpistemicView,Everaere2020BeliefMerging}.

Some precursors in the belief merging literature 
\citep{Schwind2015BeliefRevisiona,Schwind2016PromotingBeliefs,Schwind2024BeliefFlowFramework,Schwind2026CharacterizingConsensuses},
though these only look at convergence (no truth involved).

The opinion diffusion literature \citep{Grandi2015PropositionalOpinion,Christoff2017StabilityBinary,Botan2019MultiIssueOpinion,Grandi2021GamesInfluence}.

The computational literature \citep{Betz2013DebateDynamics}.

The DeGroot dynamics \citep{DeGroot1974ReachingConsensus}, and the wisdom aspect, where we have a ground truth \citep{Golub2010NaiveLearning}.


\section{Model}\label{sec:model}
% For any $n \geq 1$, $[n] = \{1, \dots, n\}$ is the set of positive integers from $1$ to $n$.


\paragraph{Agents.}
\emph{Agents} are the nodes of a finite, directed graph $G = (N, E)$, 
where $N = \{1, \dots, n\}$ is the set of agents, and $E$ is the set of edges.   
We write $N(i)$ for $i$'s out-neighborhood in $G$, 
i.e., $N(i) = \{j \mid (i, j) \in E\}$, the set of agents that $i$ has an edge to.
Intuitively, $N(i)$ is the set of agents that $i$ pays attention to $G$.  

\paragraph{Opinion Dynamics.}
There is a set $A = \{a_1, \dots a_m\}$ of \emph{atoms},
with $L$ being the set of \emph{logical formulas} constructed using these atoms
using the usual connectives ($\lnot$, $\land$, $\lor$, $\rightarrow$, $\bot$ and $\top$).
We call $2^A$ the set of \emph{possible worlds}, or \emph{truth-value assignments}.
If $\phi \in L$ is a propositional formula, $[\phi] \subseteq 2^A$ is the set of models of $\phi$,
i.e., the truth-value assignments that make $\phi$ true.

The dynamics occurs in rounds $t \in \{0, 1, \dots\}$.
For each agent $i$, $\phi_i^t \in L$ is agent $i$'s opinion at $t$.
We assume opinions at $t=0$ are given by a noise model (see below). 
After this, agents update their opinions using the opinions of their neighbors 
and a merging operator $M$:
$$
    \phi_i^{t+1} = \merge {\left(\phi_j^t \mid j \in N(i)\right)},
$$
i.e., $i$'s opinion at $t+1$ is the result of merging 
the opinions at $t$ of $i$'s neighbors in $G$.
Note that $i$'s own opinion at $t$ enters here if we assume there is an edge from $i$ to itself $G$.     

\paragraph{Merging Operators.}
What merging operators shall we use?

\paragraph{Noise Model.}
Here we describe a noise model that generates initial opinions. 

We assume, first of all, that there exists a ground truth state of the world 
$w^\star \in 2^A$.

[TBC]

\paragraph{Questions.}
Do we get that each agent's beliefs stabilize, i.e., is there some belief $\phi_i^*$ 
such that:
$$
    \lim_{t \rightarrow \infty} \phi_i^t = \phi_i^*,
$$
for any agent $i$? Under what conditions does this happen?

If agents' beliefs stabilize, do we see consensus? In other words, 
is there some belief $\phi^*$ such that $\phi_i^* = \phi^*$,
for any agent $i$? And under what conditions?

If there is consensus, do we see consensus to the truth as $n$ grows? 
Take a family of graphs $G_n$ in which the dynamics 
reaches consensus $\phi^*(n)$ as $t \rightarrow \infty$.
Call a family of graphs $G_n$ \emph{wise} if:
$$
    \lim_{n \rightarrow \infty} {\Pr}\Big[ \big[\phi^*(n)\big] = w^\star \Big] = 1,
$$   
i.e., if the consensus belief over the whole graph $G_n$ is exactly (a formula with the model) $w^\star$,
the ground truth state, with high probability, as $n$ grows.    
Our question, then, is what kind of networks are wise, and for which noise models?

\bibliography{bibliography}
% \bibliographystyle{plain}
\end{document}